% Related Work – ca. 500 Wörter
% Drei Blöcke: OER-Tools, Feature-Adoption, GitHub-Mining

Analysing feature adoption in a client-side OER tool requires drawing on
three areas of prior work:
the \emph{OER authoring tool} landscape, which defines the design space
LiaScript operates in;
\emph{feature adoption research} in educational technologies, which
establishes that gaps between designed and used features are the norm
rather than the exception;
and \emph{repository mining}, which provides the methodological foundation
for studying a tool that produces no server-side telemetry.

\subsection{OER Authoring Tools}
\label{sec:rw:tools}

Several tools address the challenge of creating interactive educational content
from plain-text sources.
H5P~\cite{Kiryakova2022,Fornasari2025} is widely adopted for creating
self-contained interactive elements (quizzes, interactive videos) that can be
embedded in any LMS, but requires a server and does not support version control.
Jupyter Book~\cite{JupyterBook2025,Betlem2025} and mdBook generate complete
websites from Markdown files and integrate naturally with Git workflows, but
are primarily oriented towards computational and technical content.
The \emph{You Only Write Thrice} approach~\cite{WriteThrice2021} demonstrates
that a single Markdown source can generate a book, a slide deck, and a
computational notebook simultaneously --- a principle LiaScript generalises
to the full OER lifecycle.
Pandoc-based workflows~\cite{Krewinkel2017} address format interoperability
but do not add interactive pedagogical elements.
Git-based OER management has also been explored with GitLab as
infrastructure~\cite{Moser2025} and with GitHub specifically for
collaborative course development~\cite{Schwiertz2024}.

LiaScript differs from all of the above in combining Markdown simplicity,
rich interactive affordances, and a fully client-side execution model
that requires neither a server nor a pre-installed runtime~\cite{Baur2025}.

\subsection{Feature Adoption in Educational Technologies}
\label{sec:rw:adoption}

A recurring finding in educational technology research is that the number of
available features in a system far exceeds the number of features that are
actually used.
Studies on learning management systems document this pattern consistently:
Hijazi et al.~\cite{Hijazi2020} found a strong positive correlation between
feature \emph{awareness} and feature \emph{usage} in Moodle, with advanced
features such as SCORM packages and workshop modules remaining largely unknown
and therefore unused.
Awad et al.~\cite{Awad2019} observed similarly uneven utilisation of LMS
features even at well-resourced institutions.

The Technology Acceptance Model (TAM) and its extensions provide a theoretical
framework for these observations~\cite{Granic2022}.
Feng et al.~\cite{Feng2025} identify \emph{Effort Expectancy} --- the perceived
complexity of using a feature --- as the dominant predictor of non-adoption in
higher education contexts, outweighing perceived usefulness for features that
require significant authoring effort.

Crucially, all of these studies were conducted on server-based systems that
generate usage logs.
No prior work analyses feature adoption in a decentralised, client-side OER
tool where no telemetry is available and repository content is the only
observable proxy for usage behaviour.

\subsection{Mining Educational Repositories}
\label{sec:rw:mining}

Mining software repositories (MSR) is an established methodology for
understanding how tools are used in practice~\cite{Cosentino2017,Cosentino2016}.
Schröder and Pfänder~\cite{Schroeder2020} specifically explore GitHub as a
platform for OER, and Santos-Hermosa et al.~\cite{SantosHermosa2017} analyse
reuse indicators across OER repositories.
Both studies confirm that repository metadata --- commit history, contributor
counts, file content --- provides meaningful signals about authoring behaviour,
while also cautioning about the incompleteness of public repository data as a
usage proxy.

Our work extends this line of research by applying corpus analysis to a
specific OER authoring language, treating the GitHub corpus as the \emph{only
available substitute} for telemetry data in a privacy-preserving,
client-side tool.
